\section{Evidencias del uso de herramientas de apoyo}

En este apartado se incluirán capturas de pantalla que permitan identificar las consultas realizadas a las herramientas de apoyo durante el desarrollo de la actividad.

Las evidencias se presentarán mediante figuras numeradas, cada una con un título que describa el contenido de la captura. En el texto se hará referencia a las figuras correspondientes: la Figura~\ref{fig:clasificacion} corresponde a la consulta sobre la clasificación de requisitos, la Figura~\ref{fig:reformulacion} a la consulta sobre la reformulación de requisitos de calidad y la Figura~\ref{fig:iso25010} a la consulta sobre las características de calidad de ISO/IEC 25010:2023.

\begin{figure}[H]
  \centering
  \capturaimagen{captura1}{Captura de pantalla 1\\[0.4em]Insertar aquí la captura real de la consulta relacionada con la clasificación de los 20 requisitos.\\[0.4em]Archivo esperado: imagenes/captura1.png}
  \caption{Consulta de apoyo sobre la redacción, ortografía y gramática del trabajo.}
  \label{fig:clasificacion}
\end{figure}


\Needspace{9\baselineskip}
\subsection{Registro de herramientas utilizadas}

\begingroup
\small
\setlength{\tw}{\dimexpr\linewidth-8\tabcolsep-5\arrayrulewidth\relax}
\begin{longtable}{|L{0.17\tw}|L{0.10\tw}|L{0.38\tw}|L{0.35\tw}|}
\caption{Registro de herramientas utilizadas.}\label{tab:herramientas}\\
\hline
\filaencabezado \encabezado{Herramienta} & \encabezado{Versión} & \encabezado{Propósito de uso} & \encabezado{Verificación} \\ \hline
\endfirsthead
\hline
\filaencabezado \encabezado{Herramienta} & \encabezado{Versión} & \encabezado{Propósito de uso} & \encabezado{Verificación} \\ \hline
\endhead
ChatGPT & GPT-6 & Solicitar una revisión exhaustiva de la redacción, ortografía, gramática y coherencia del trabajo académico, manteniendo su estructura y contenido original. & Se revisaron y corrigieron errores de escritura, concordancia y claridad. Se conservaron las clasificaciones, métricas, referencias y conclusiones, y se señalaron los aspectos técnicos que requieren validación. \\ \hline
\end{longtable}
\endgroup
